\par\addvspace{2.4ex}\noindent
\begin{minipage}{\linewidth}
\qhead{1}{1.1}{向量计算}{ch01:sec:1.1}

% ---- item_id=e0a6a6b00183 seq=1.1-001 ----
\begin{q}{1.1-001}{1995数一}{ch01:q:e0a6a6b00183}
设 $(\mathbf{a}\times\mathbf{b})\cdot\mathbf{c}=2$，则 $[(\mathbf{a}+\mathbf{b})\times(\mathbf{b}+\mathbf{c})]\cdot(\mathbf{c}+\mathbf{a})=\underline{\hspace{4em}}$
\end{q}
\end{minipage}
\markboth{1.1\quad 向量计算}{1.1\quad 向量计算}


% ---- item_id=6d50cbd015f3 seq=1.1-002 ----
\begin{q}{1.1-002}{880第四章基础填空1}{ch01:q:6d50cbd015f3}
设向量 $a=(-1,3,0), b=(3,1,0), |c|=r$(常数). 当 $c$ 满足 $a=b\times c$ 时, $r$ 的最小值为 \underline{\hspace{4em}} .
\end{q}

% ---- item_id=7c0239494ee8 seq=1.1-003 ----
\begin{q}{1.1-003}{880第四章基础填空2}{ch01:q:7c0239494ee8}
设向量 $a=(2,-3,1),b=(1,-2,3),c=(2,1,2)$，向量 $r$ 满足 $r\perp a,r\perp b,\mathrm{Prj}_{c}r=14$，则 $r=\underline{\hspace{4em}}$.
\end{q}

% ---- item_id=4ed3597e1bf4 seq=1.1-004 ----
\begin{q}{1.1-004}{880第四章基础选择1}{ch01:q:4ed3597e1bf4}
(1) 设向量 $a=(1,2,1),b=(-1,0,2),c=(0,k,-3)$ 共面，则 $k=(\quad)$．
\begin{opts}{0.235}
\opt{A}{$1$}\hfill\opt{B}{$2$}\hfill\opt{C}{$-1$}\hfill\opt{D}{$-2$}
\end{opts}
\end{q}

% ---- item_id=351c298870a7 seq=1.1-005 ----
\begin{q}{1.1-005}{880第四章基础填空11}{ch01:q:351c298870a7}
设 $\alpha$ 与 $\beta$ 均为单位向量，其夹角为 $\frac{\pi}{6}$，则以 $\alpha+2\beta$ 与 $3\alpha+\beta$ 为邻边的平行四边形的面积为 $\underline{\hspace{4em}}$．
\end{q}

% ---- item_id=2d0194f9e3d7 seq=1.1-006 ----
\begin{q}{1.1-006}{880第四章基础填空12}{ch01:q:2d0194f9e3d7}
设 $\alpha$ 与 $\beta$ 是非零常向量，$|\beta|=2$（$|\beta|$ 表示 $\beta$ 的模），$\alpha$ 与 $\beta$ 之间的夹角为 $\frac{\pi}{3}$，求 $\displaystyle \lim\limits_{x\to 0}\frac{|\alpha+x\beta|-|\alpha|}{x}=$
\end{q}

% ---- item_id=5e228db2d0e8 seq=1.1-007 ----
\begin{q}{1.1-007}{900第四章A类4}{ch01:q:5e228db2d0e8}
若向量 $\boldsymbol{\alpha}=(1,0,0)$，$\boldsymbol{\beta}=(0,1,0)$，$\boldsymbol{\gamma}=(0,0,1)$，则 $[(\boldsymbol{\alpha}+\boldsymbol{\beta}+\boldsymbol{\gamma})\times(\boldsymbol{\alpha}+\boldsymbol{\beta})]\cdot\boldsymbol{\alpha}=\underline{\hspace{4em}}$。
\end{q}

% ---- item_id=c0a182647e09 seq=1.1-008 ----
\begin{q}{1.1-008}{660数一137}{ch01:q:c0a182647e09}
已知三个向量 $a,b,c$，其中 $c\perp a,c\perp b$，$a$ 与 $b$ 的夹角为 $\frac{\pi}{6}$，$|a|=6,|b|=|c|=3$，则 $(a\times b)\cdot c=\underline{\hspace{4em}}$
\end{q}

\par\addvspace{2.4ex}\noindent
\begin{minipage}{\linewidth}
\qhead{1}{1.2}{空间直线与平面}{ch01:sec:1.2}

\qhead{2}{1.2.1}{求直线方程}{ch01:sec:1.2.1}

\qhead{3}{1.2.1.1}{已知线面关系}{ch01:sec:1.2.1.1}

% ---- item_id=90dd7792133f seq=1.2.1.1-001 ----
\begin{q}{1.2.1.1-001}{880第四章基础解答5}{ch01:q:90dd7792133f}
求过点 $(-1,2,3)$，垂直于直线 $\frac{x}{4}=\frac{y}{5}=\frac{z}{6}$ 且平行于平面 $7x+8y+9z+10=0$ 的直线方程.
\end{q}
\end{minipage}
\markboth{1.2.1.1\quad 已知线面关系}{1.2.1.1\quad 已知线面关系}


% ---- item_id=5c10d1bdeb0c seq=1.2.1.1-002 ----
\begin{q}{1.2.1.1-002}{660数一139}{ch01:q:5c10d1bdeb0c}
过点 $P(-1,0,4)$ 且与平面 $3x-4y+z+10=0$ 平行，又与直线 $L:\frac{x+1}{1}=\frac{y-3}{1}=\frac{z}{2}$ 相交的直线方程是 $\underline{\hspace{4em}}$．
\end{q}

\par\addvspace{2.4ex}\noindent
\begin{minipage}{\linewidth}
\qhead{3}{1.2.1.2}{公垂线}{ch01:sec:1.2.1.2}

% ---- item_id=b9fcfcbe4a10 seq=1.2.1.2-001 ----
\begin{q}{1.2.1.2-001}{880第四章基础解答6}{ch01:q:b9fcfcbe4a10}
求与直线 $L_1:x+2=3-y=z+1$ 和 $L_2:\frac{x+4}{2}=y=\frac{z-4}{3}$ 都垂直相交的直线方程.
\end{q}
\end{minipage}
\markboth{1.2.1.2\quad 公垂线}{1.2.1.2\quad 公垂线}


% ---- item_id=6a1a2ba39c32 seq=1.2.1.2-002 ----
\begin{q}{1.2.1.2-002}{880第四章基础解答7}{ch01:q:6a1a2ba39c32}
求直线 $L_1:\frac{x-3}{2}=y=\frac{z-1}{0}$ 与 $L_2:\frac{x+1}{1}=\frac{y-2}{0}=z$ 的公垂线方程.
\end{q}

\par\addvspace{2.4ex}\noindent
\begin{minipage}{\linewidth}
\qhead{3}{1.2.1.3}{投影线}{ch01:sec:1.2.1.3}

% ---- item_id=8922c4948527 seq=1.2.1.3-001 ----
\begin{q}{1.2.1.3-001}{880第四章基础解答4}{ch01:q:8922c4948527}
求直线 $L:\frac{x+2}{3}=\frac{2-y}{1}=\frac{z+1}{2}$ 在平面 $\pi:2x+3y+3z-8=0$ 上的投影直线方程.
\end{q}
\end{minipage}
\markboth{1.2.1.3\quad 投影线}{1.2.1.3\quad 投影线}


% ---- item_id=5a893c936e10 seq=1.2.1.3-002 ----
\begin{q}{1.2.1.3-002}{880第四章基础填空7}{ch01:q:5a893c936e10}
直线 $L:\begin{cases} x+2y-3z=2,\\ 2x-y+z=3 \end{cases}$，在平面 $z=0$ 上的投影为 $\underline{\hspace{4em}}$，在平面 $z=1$ 上的投影为 $\underline{\hspace{4em}}$.
\end{q}

\par\addvspace{2.4ex}\noindent
\begin{minipage}{\linewidth}
\qhead{2}{1.2.2}{求平面方程}{ch01:sec:1.2.2}

\qhead{3}{1.2.2.1}{线线关系}{ch01:sec:1.2.2.1}

% ---- item_id=0fd76c2fb7a7 seq=1.2.2.1-001 ----
\begin{q}{1.2.2.1-001}{1987数一}{ch01:q:0fd76c2fb7a7}
与两直线 \[\begin{cases} x=1, \\ y=-1+t, \\ z=2+t \end{cases}\] 及 $\displaystyle \frac{x+1}{1}=\frac{y+2}{2}=\frac{z-1}{1}$ 都平行，且过原点的平面方程为 \underline{\hspace{4em}}
\end{q}
\end{minipage}
\markboth{1.2.2.1\quad 线线关系}{1.2.2.1\quad 线线关系}


% ---- item_id=917cd7a357c9 seq=1.2.2.1-002 ----
\begin{q}{1.2.2.1-002}{1991数一}{ch01:q:917cd7a357c9}
（3）已知两条直线 $L_1:\frac{x-1}{1}=\frac{y-2}{0}=\frac{z-3}{-1},L_2:\frac{x+2}{2}=\frac{y-1}{1}=\frac{z}{1}$，则过 $L_1$ 且平行于 $L_2$ 的平面方程是\underline{\hspace{4em}}
\end{q}

% ---- item_id=84137a6bcad6 seq=1.2.2.1-003 ----
\begin{q}{1.2.2.1-003}{900第四章A类6}{ch01:q:84137a6bcad6}
若直线 $l_1:\frac{x-4}{1}=\frac{y-2}{3}=\frac{z-1}{5}$ 在平面 $\pi$ 上的投影直线为 $l_2:\frac{x-4}{3}=\frac{y-2}{2}=\frac{z-1}{1}$，则平面 $\pi$ 的方程为 $\underline{\hspace{4em}}$.
\end{q}

% ---- item_id=9163b7653010 seq=1.2.2.1-004 ----
\begin{q}{1.2.2.1-004}{900第四章A类7}{ch01:q:9163b7653010}
10 与两直线 $\begin{cases} x=2,\\ y=t,\\ z=1+t \end{cases}$ 及 $\frac{x}{1}=\frac{y+1}{2}=\frac{z-1}{3}$ 都平行，且过原点的平面方程为 $\underline{\hspace{4em}}$.
\end{q}

% ---- item_id=c67ba335bb22 seq=1.2.2.1-005 ----
\begin{q}{1.2.2.1-005}{900第四章A类8}{ch01:q:c67ba335bb22}
11 已知两条直线的方程是
$\displaystyle L_1:\frac{x-1}{1}=\frac{y-2}{1}=\frac{z-1}{-1},\quad L_2:\frac{x+2}{2}=\frac{y-1}{1}=\frac{z}{1},$
则过 $L_1$ 且平行于 $L_2$ 的平面方程是 \underline{\hspace{4em}}.
\end{q}

\par\addvspace{2.4ex}\noindent
\begin{minipage}{\linewidth}
\qhead{3}{1.2.2.2}{点线关系}{ch01:sec:1.2.2.2}

% ---- item_id=f35c05ce13e6 seq=1.2.2.2-001 ----
\begin{q}{1.2.2.2-001}{1990数一}{ch01:q:f35c05ce13e6}
过点 $M(1,2,-1)$ 且与直线 \[\begin{cases} x=-t+2 \\ y=3t-4 \\ z=t-1 \end{cases}\] 垂直的平面方程是 \underline{\hspace{4em}}
\end{q}
\end{minipage}
\markboth{1.2.2.2\quad 点线关系}{1.2.2.2\quad 点线关系}


% ---- item_id=f2ac298199cb seq=1.2.2.2-002 ----
\begin{q}{1.2.2.2-002}{880第四章基础填空3}{ch01:q:f2ac298199cb}
过点 $(2,0,-3)$ 且与直线 $\begin{cases} x-2y+4z-7=0,\\ 3x+5y-2z+1=0 \end{cases}$ 垂直的平面方程为 $\underline{\hspace{4em}}$.
\end{q}

% ---- item_id=4bedabdabb68 seq=1.2.2.2-003 ----
\begin{q}{1.2.2.2-003}{880第四章基础填空8}{ch01:q:4bedabdabb68}
过点 $A(1,1,-1),B(-2,-2,2)$ 和 $C(1,-1,2)$ 三点的平面方程为 \underline{\hspace{4em}}.
\end{q}

\par\addvspace{2.4ex}\noindent
\begin{minipage}{\linewidth}
\qhead{3}{1.2.2.3}{平面束}{ch01:sec:1.2.2.3}

% ---- item_id=5ec9aa5b9008 seq=1.2.2.3-001 ----
\begin{q}{1.2.2.3-001}{880第四章基础解答2}{ch01:q:5ec9aa5b9008}
设平面 $\pi$ 与点 $P(1,2,1)$ 的距离为 $1$，且过直线 \[L:\begin{cases}3x-2y+2=0,\\ x-2y-z+6=0,\end{cases}\] 求平面 $\pi$ 的方程.
\end{q}
\end{minipage}
\markboth{1.2.2.3\quad 平面束}{1.2.2.3\quad 平面束}


% ---- item_id=9bcde0ba1deb seq=1.2.2.3-002 ----
\begin{q}{1.2.2.3-002}{880第四章基础解答3}{ch01:q:9bcde0ba1deb}
设平面 $\pi$ 过直线 $L:\begin{cases} x+5y+z=0, \\ x-z+4=0, \end{cases}$ 且与平面 $\pi_1:x-4y-8z+12=0$ 的夹角为 $45^\circ$，求平面 $\pi$ 的方程.
\end{q}

\par\addvspace{2.4ex}\noindent
\begin{minipage}{\linewidth}
\qhead{3}{1.2.2.4}{点面关系}{ch01:sec:1.2.2.4}

% ---- item_id=6c8305891950 seq=1.2.2.4-001 ----
\begin{q}{1.2.2.4-001}{1996数一}{ch01:q:6c8305891950}
设一平面经过原点及 $(6,-3,2)$，且与平面 $4x-y+2z=8$ 垂直，则此平面方程为 $\underline{\hspace{4em}}$
\end{q}
\end{minipage}
\markboth{1.2.2.4\quad 点面关系}{1.2.2.4\quad 点面关系}


\par\addvspace{2.4ex}\noindent
\begin{minipage}{\linewidth}
\qhead{3}{1.2.2.5}{面面关系}{ch01:sec:1.2.2.5}

% ---- item_id=de10b182a010 seq=1.2.2.5-001 ----
\begin{q}{1.2.2.5-001}{880第四章基础解答1}{ch01:q:de10b182a010}
求平行于平面 $x+y+z=9$ 且与球面 $x^2+y^2+z^2=4$ 相切的平面方程。
\end{q}
\end{minipage}
\markboth{1.2.2.5\quad 面面关系}{1.2.2.5\quad 面面关系}


\par\addvspace{2.4ex}\noindent
\begin{minipage}{\linewidth}
\qhead{3}{1.2.2.6}{面距关系}{ch01:sec:1.2.2.6}

% ---- item_id=f030f6142839 seq=1.2.2.6-001 ----
\begin{q}{1.2.2.6-001}{660数一136}{ch01:q:f030f6142839}
平行于平面 $5x-14y+2z+36=0$ 且距此平面距离为 $3$ 的平面方程为 $\underline{\hspace{4em}}$．
\end{q}
\end{minipage}
\markboth{1.2.2.6\quad 面距关系}{1.2.2.6\quad 面距关系}


\par\addvspace{2.4ex}\noindent
\begin{minipage}{\linewidth}
\qhead{2}{1.2.3}{求夹角}{ch01:sec:1.2.3}

% ---- item_id=c29b9607d5ab seq=1.2.3-001 ----
\begin{q}{1.2.3-001}{1993数一; 880第四章基础选择2}{ch01:q:c29b9607d5ab}
(3) 设有直线 $L_1:\frac{x-1}{1}=\frac{y-5}{-2}=\frac{z+8}{1}$ 与 $L_2:\begin{cases} x-y=6,\\ 2y+z=3, \end{cases}$ 则 $L_1$ 与 $L_2$ 的夹角为（\quad）
\begin{opts}{0.235}
\opt{A}{$\frac{\pi}{6}.$}\hfill\opt{B}{$\frac{\pi}{4}.$}\hfill\opt{C}{$\frac{\pi}{3}.$}\hfill\opt{D}{$\frac{\pi}{2}.$}
\end{opts}
\end{q}
\end{minipage}
\markboth{1.2.3\quad 求夹角}{1.2.3\quad 求夹角}


% ---- item_id=2ca5e1ed5a7a seq=1.2.3-002 ----
\begin{q}{1.2.3-002}{900第四章A类1}{ch01:q:2ca5e1ed5a7a}
已知直线 $l_1$ 过点 $(1,2,3)$ 和 $(9,4,1)$，则直线 $l_1$ 与直线 $l_2$：$\begin{cases} x+2y+3z=1,\\ 9x+4y+z=1 \end{cases}$ 的夹角为（  ）
\begin{opts}{0.235}
\opt{A}{$0.$}\hfill\opt{B}{$\frac{\pi}{6}.$}\hfill\opt{C}{$\frac{\pi}{3}.$}\hfill\opt{D}{$\frac{\pi}{2}.$}
\end{opts}
\end{q}

% ---- item_id=fc3d0c0ae542 seq=1.2.3-003 ----
\begin{q}{1.2.3-003}{660数一341}{ch01:q:fc3d0c0ae542}
341 直线 $L:\frac{x-2}{2}=\frac{y-1}{1}=\frac{z-3}{1}$ 与平面 $\Pi:x-y+2z+4=0$ 的夹角为
\begin{opts}{0.235}
\opt{A}{$\pi$.}\hfill\opt{B}{$\frac{\pi}{3}$.}\hfill\opt{C}{$\frac{\pi}{6}$.}\hfill\opt{D}{$\frac{\pi}{2}$.}
\end{opts}
\end{q}

\par\addvspace{2.4ex}\noindent
\begin{minipage}{\linewidth}
\qhead{2}{1.2.4}{求距离}{ch01:sec:1.2.4}

% ---- item_id=6e2e4e0c3b10 seq=1.2.4-001 ----
\begin{q}{1.2.4-001}{2006数一}{ch01:q:6e2e4e0c3b10}
点 $(2,1,0)$ 到平面 $3x+4y+5z=0$ 的距离 $d = \underline{\hspace{4em}}$
\end{q}
\end{minipage}
\markboth{1.2.4\quad 求距离}{1.2.4\quad 求距离}


% ---- item_id=288aa8438128 seq=1.2.4-002 ----
\begin{q}{1.2.4-002}{880第四章基础填空4}{ch01:q:288aa8438128}
点 $P(1,-1,2)$ 到平面 $\pi: 2x-y+5z-12=0$ 的距离 $d=\underline{\hspace{4em}}$。
\end{q}

% ---- item_id=90d1665c2e23 seq=1.2.4-003 ----
\begin{q}{1.2.4-003}{880第四章基础填空5}{ch01:q:90d1665c2e23}
点 $P(1,-1,0)$ 到直线 $L:\frac{x}{3}=\frac{y}{3}=\frac{z+1}{2}$ 的距离 $d=\underline{\hspace{4em}}$.
\end{q}

% ---- item_id=61bbfba088f3 seq=1.2.4-004 ----
\begin{q}{1.2.4-004}{660数一138}{ch01:q:61bbfba088f3}
点 $M(3,2,6)$ 到直线 $\frac{x}{1}=\frac{y+7}{2}=\frac{z-3}{-1}$ 的距离为 \underline{\hspace{4em}}.
\end{q}

% ---- item_id=f00c77c8a10f seq=1.2.4-005 ----
\begin{q}{1.2.4-005}{880第四章基础填空6}{ch01:q:f00c77c8a10f}
直线 $L_1:\begin{cases} x-1=0, \\ y=z \end{cases}$ 与 $L_2:\begin{cases} x+2y=0, \\ z+2=0 \end{cases}$ 之间的距离 $d=\underline{\hspace{4em}}$ .
\end{q}

% ---- item_id=4d9b6c54c6b5 seq=1.2.4-006 ----
\begin{q}{1.2.4-006}{900第四章A类5}{ch01:q:4d9b6c54c6b5}
点 $(1,2,0)$ 到平面 $3x+4y+5z=0$ 的距离 $d=\underline{\hspace{4em}}$.
\end{q}

% ---- item_id=b27d25eb12e6 seq=1.2.4-007 ----
\begin{q}{1.2.4-007}{1000第17章强化2}{ch01:q:b27d25eb12e6}
求曲面 $4z=3x^2-2xy+3y^2$ 上的点到平面 $x+y-4z=1$ 的最短距离.
\end{q}

\par\addvspace{2.4ex}\noindent
\begin{minipage}{\linewidth}
\qhead{2}{1.2.5}{直线平面的位置关系}{ch01:sec:1.2.5}

% ---- item_id=54aed8c6af29 seq=1.2.5-001 ----
\begin{q}{1.2.5-001}{1995数一; 1000第17章基础2; 660数一342; 880第四章基础选择3}{ch01:q:54aed8c6af29}
（1）设有直线 $L$：\[\begin{cases} x+3y+2z+1=0,\\ 2x-y-10z+3=0 \end{cases}\] 及平面 $\pi:4x-2y+z-2=0$，则直线 $L$（\quad）
\begin{opts}{0.235}
\opt{A}{平行于 $\pi$．}\hfill\opt{B}{在 $\pi$ 上．}\hfill\opt{C}{垂直于 $\pi$．}\hfill\opt{D}{与 $\pi$ 斜交．}
\end{opts}
\end{q}
\end{minipage}
\markboth{1.2.5\quad 直线平面的位置关系}{1.2.5\quad 直线平面的位置关系}


% ---- item_id=ad9ac7921593 seq=1.2.5-002 ----
\begin{q}{1.2.5-002}{900第四章A类2}{ch01:q:ad9ac7921593}
已知平面 $\pi$ 过原点且直线 $l:\frac{x-3}{1}=\frac{y+2}{2}=\frac{z-6}{3}$ 在平面 $\pi$ 上，则下列直线在平面 $\pi$ 上的是（ ）
\begin{opts}{0.48}
\opt{A}{$l_1:\frac{x+1}{1}=\frac{y-6}{-6}=\frac{z}{-1}$.}\hfill\opt{B}{$l_2:\frac{x-2}{2}=\frac{y+4}{-3}=\frac{z-3}{3}$.}\\[0.3ex]
\opt{C}{$l_3:\frac{x-1}{2}=\frac{y-1}{-4}=\frac{z-3}{3}$.}\hfill\opt{D}{$l_4:\frac{x+1}{1}=\frac{y-6}{2}=\frac{z}{3}$.}
\end{opts}
\end{q}

\par\addvspace{2.4ex}\noindent
\begin{minipage}{\linewidth}
\qhead{1}{1.3}{空间曲线与曲面}{ch01:sec:1.3}

\qhead{2}{1.3.1}{空间曲线}{ch01:sec:1.3.1}

\qhead{3}{1.3.1.1}{求曲线方程}{ch01:sec:1.3.1.1}

% ---- item_id=3af2d1af57fd seq=1.3.1.1-001 ----
\begin{q}{1.3.1.1-001}{880第四章基础填空10}{ch01:q:3af2d1af57fd}
曲线 $\begin{cases} z=2x^2+3y^2, \\ z=12-x^2-3y^2 \end{cases}$ 在 $xOy$ 面上的投影曲线方程为 $\underline{\hspace{4em}}$.
\end{q}
\end{minipage}
\markboth{1.3.1.1\quad 求曲线方程}{1.3.1.1\quad 求曲线方程}


\par\addvspace{2.4ex}\noindent
\begin{minipage}{\linewidth}
\qhead{3}{1.3.1.2}{切线与法平面}{ch01:sec:1.3.1.2}

% ---- item_id=e8daefd3df6f seq=1.3.1.2-001 ----
\begin{q}{1.3.1.2-001}{1992数一}{ch01:q:e8daefd3df6f}
（3）在曲线 $x=t,y=-t^2,z=t^3$ 的所有切线中,与平面 $x+2y+z=4$ 平行的切线（ ）
\begin{opts}{0.235}
\opt{A}{只有 1 条.}\hfill\opt{B}{只有 2 条.}\hfill\opt{C}{至少有 3 条.}\hfill\opt{D}{不存在.}
\end{opts}
\end{q}
\end{minipage}
\markboth{1.3.1.2\quad 切线与法平面}{1.3.1.2\quad 切线与法平面}


% ---- item_id=149ea2168dbf seq=1.3.1.2-002 ----
\begin{q}{1.3.1.2-002}{2001数一}{ch01:q:149ea2168dbf}
（2）设函数 $f(x,y)$ 在点 $(0,0)$ 附近有定义，且 $f_x'(0,0)=3$，$f_y'(0,0)=1$，则（\quad）
\begin{opts}{0.97}
\opt{A}{$\left.dz\right|_{(0,0)}=3dx+dy$.}\\[0.3ex]
\opt{B}{曲面 $z=f(x,y)$ 在点 $(0,0,f(0,0))$ 的法向量为 $(3,1,1)$.}\\[0.3ex]
\opt{C}{曲线 $\begin{cases} z=f(x,y),\\ y=0 \end{cases}$ 在点 $(0,0,f(0,0))$ 的切向量为 $(1,0,3)$.}\\[0.3ex]
\opt{D}{曲线 $\begin{cases} z=f(x,y),\\ y=0 \end{cases}$ 在点 $(0,0,f(0,0))$ 的切向量为 $(3,0,1)$.}
\end{opts}
\end{q}

% ---- item_id=06730c0f0539 seq=1.3.1.2-003 ----
\begin{q}{1.3.1.2-003}{1000第17章强化3}{ch01:q:06730c0f0539}
空间曲线 $L$: \[ \begin{cases} y^2=z, \\ x=2(y-1) \end{cases} \] 在 $y=1$ 处的切线方程为 \underline{\hspace{4em}}。
\end{q}

% ---- item_id=31792a89007c seq=1.3.1.2-004 ----
\begin{q}{1.3.1.2-004}{660数一343}{ch01:q:31792a89007c}
曲线 \[\begin{cases} x^2+y^2+z^2=2 \\ x+y+z=0 \end{cases}\] 在点 $(1,-1,0)$ 处的切线方程为
\begin{opts}{0.48}
\opt{A}{$\frac{x-1}{2}=y+1=z.$}\hfill\opt{B}{$\frac{x-1}{2}=\frac{y+1}{2}=\frac{z}{3}.$}\\[0.3ex]
\opt{C}{$\frac{x-1}{-1}=\frac{y+1}{-1}=\frac{z}{1}.$}\hfill\opt{D}{$x-1=y+1=-\frac{z}{2}.$}
\end{opts}
\end{q}

% ---- item_id=dfddb7049086 seq=1.3.1.2-005 ----
\begin{q}{1.3.1.2-005}{1000第17章强化4}{ch01:q:dfddb7049086}
曲线 $\begin{cases} x^2+2y^2+3z^2=6, \\ x+y+z=3 \end{cases}$ 在点 $(1,1,1)$ 处的切线方程为 $\underline{\hspace{4em}}$．
\end{q}

% ---- item_id=1b611462499d seq=1.3.1.2-006 ----
\begin{q}{1.3.1.2-006}{1000第17章强化6}{ch01:q:1b611462499d}
设函数 $z=f(x,y)$ 在点 $(0,0)$ 附近有定义，且 $f'_x(0,0)=3$，则曲线
\[\begin{cases}
z=f(x,y),\\
y=0
\end{cases}\]
在点 $(0,0,f(0,0))$ 处的法平面方程为 $\underline{\hspace{4em}}$.
\end{q}

\par\addvspace{2.4ex}\noindent
\begin{minipage}{\linewidth}
\qhead{2}{1.3.2}{空间曲面}{ch01:sec:1.3.2}

\qhead{3}{1.3.2.1}{旋转曲面方程}{ch01:sec:1.3.2.1}

% ---- item_id=546d26dfe9d9 seq=1.3.2.1-001 ----
\begin{q}{1.3.2.1-001}{1998数一}{ch01:q:546d26dfe9d9}
求直线 $l:\frac{x-1}{1}=\frac{y}{1}=\frac{z-1}{-1}$ 在平面 $\pi:x-y+2z-1=0$ 上的投影直线 $l_0$ 的方程，并求 $l_0$ 绕 $y$ 轴旋转一周所成曲面的方程.
\end{q}
\end{minipage}
\markboth{1.3.2.1\quad 旋转曲面方程}{1.3.2.1\quad 旋转曲面方程}


% ---- item_id=028971b57920 seq=1.3.2.1-002 ----
\begin{q}{1.3.2.1-002}{880第四章基础填空9}{ch01:q:028971b57920}
曲线 $\begin{cases} x^2+2z^2=4, \\ y=0 \end{cases}$ 绕 $z$ 轴旋转一周所得的旋转曲面方程为 \underline{\hspace{4em}}.
\end{q}

% ---- item_id=87c152010633 seq=1.3.2.1-003 ----
\begin{q}{1.3.2.1-003}{880第四章基础解答8}{ch01:q:87c152010633}
求直线 $\frac{x-1}{0}=\frac{y}{1}=\frac{z-1}{2}$ 绕 $z$ 轴旋转一周所得的旋转曲面方程.
\end{q}

% ---- item_id=033407f4aec1 seq=1.3.2.1-004 ----
\begin{q}{1.3.2.1-004}{880第四章基础解答9}{ch01:q:033407f4aec1}
求直线 $L:\frac{x-1}{3}=\frac{y-2}{4}=\frac{z+1}{1}$ 绕直线 $\begin{cases} x=2, \\ y=3 \end{cases}$ 旋转一周所得的曲面方程.
\end{q}

% ---- item_id=9c2eaef66731 seq=1.3.2.1-005 ----
\begin{q}{1.3.2.1-005}{900第四章A类9}{ch01:q:9c2eaef66731}
12 求直线 $l:\frac{x}{2}=\frac{y-1}{2}=\frac{z-1}{-1}$ 在平面 $\pi:x-y+2z-1=0$ 上的投影直线 $l_0$ 的方程，并求 $l_0$ 绕 $y$ 轴旋转一周所成曲面的方程.
\end{q}

% ---- item_id=0eab31168444 seq=1.3.2.1-006 ----
\begin{q}{1.3.2.1-006}{900第四章A类3}{ch01:q:0eab31168444}
6 直线 $l:\begin{cases}3x+y+2z-2=0,\\x+2y-z-\frac{3}{2}=0\end{cases}$ 绕 $z$ 轴旋转一周的曲面记为 $\Sigma$，则下列说法中，正确的是（ ）
\begin{opts}{0.48}
\opt{A}{直线 $x-\frac{1}{2}=\frac{1}{2}-y=-z$ 在曲面 $\Sigma$ 上.}\hfill\opt{B}{直线 $\begin{cases}y=\sqrt{2}z,\\x=0\end{cases}$ 在 $\Sigma$ 上.}\\[0.3ex]
\opt{C}{曲面 $\Sigma$ 的方程为 $-x^2+y^2+2z^2=\frac{1}{2}$.}\hfill\opt{D}{曲面 $\Sigma$ 的方程为 $x^2-y^2+2z^2=\frac{1}{2}$.}
\end{opts}
\end{q}

% ---- item_id=72f7e6a3d850 seq=1.3.2.1-007 ----
\begin{q}{1.3.2.1-007}{2009数一}{ch01:q:72f7e6a3d850}
椭球面 $S_1$ 是椭圆 $\frac{x^2}{4}+\frac{y^2}{3}=1$ 绕 $x$ 轴旋转而成，圆锥面 $S_2$ 是由过点 $(4,0)$ 且与椭圆 $\frac{x^2}{4}+\frac{y^2}{3}=1$ 相切的直线绕 $x$ 轴旋转而成。

（I）求 $S_1$ 及 $S_2$ 的方程；

（II）求 $S_1$ 与 $S_2$ 之间的立体的体积。
\end{q}

\par\addvspace{2.4ex}\noindent
\begin{minipage}{\linewidth}
\qhead{3}{1.3.2.2}{切平面与法线}{ch01:sec:1.3.2.2}

% ---- item_id=88028149f8a4 seq=1.3.2.2-001 ----
\begin{q}{1.3.2.2-001}{1994数一}{ch01:q:88028149f8a4}
曲面 $z-\mathrm{e}^z+2xy=3$ 在点 $(1,2,0)$ 处的切平面方程为 $\underline{\hspace{4em}}$
\end{q}
\end{minipage}
\markboth{1.3.2.2\quad 切平面与法线}{1.3.2.2\quad 切平面与法线}


% ---- item_id=c4aa6c6a2b33 seq=1.3.2.2-002 ----
\begin{q}{1.3.2.2-002}{1000第17章强化7}{ch01:q:c4aa6c6a2b33}
曲面 $e^z-xz+y=3$ 在点 $(1,2,0)$ 处的切平面方程为 \underline{\hspace{4em}}.
\end{q}

% ---- item_id=fe764c1fc1cc seq=1.3.2.2-003 ----
\begin{q}{1.3.2.2-003}{2003数一}{ch01:q:fe764c1fc1cc}
曲面 $z=x^{2}+y^{2}$ 与平面 $2x+4y-z=0$ 平行的切平面的方程是 $\underline{\hspace{4em}}$
\end{q}

% ---- item_id=e154baf95b1f seq=1.3.2.2-004 ----
\begin{q}{1.3.2.2-004}{2013数一}{ch01:q:e154baf95b1f}
(2) 曲面 $x^2+\cos(xy)+yz+x=0$ 在点 $(0,1,-1)$ 处的切平面方程为（ ）
\begin{opts}{0.235}
\opt{A}{$x-y+z=-2$}\hfill\opt{B}{$x+y+z=0$}\hfill\opt{C}{$x-2y+z=-3$}\hfill\opt{D}{$x-y-z=0$}
\end{opts}
\end{q}

% ---- item_id=2d57d3095689 seq=1.3.2.2-005 ----
\begin{q}{1.3.2.2-005}{2014数一}{ch01:q:2d57d3095689}
曲面 $z=x^2(1-\sin y)+y^2(1-\sin x)$ 在点 $(1,0,1)$ 处的切平面方程为 \underline{\hspace{4em}}
\end{q}

% ---- item_id=4240020772f4 seq=1.3.2.2-006 ----
\begin{q}{1.3.2.2-006}{2018数一}{ch01:q:4240020772f4}
(2) 过点 $(1,0,0),(0,1,0)$，且与曲面 $z=x^2+y^2$ 相切的平面为（　　）
\begin{opts}{0.48}
\opt{A}{$z=0$ 与 $x+y-z=1$.}\hfill\opt{B}{$z=0$ 与 $2x+2y-z=2$.}\\[0.3ex]
\opt{C}{$x=y$ 与 $x+y-z=1$.}\hfill\opt{D}{$x=y$ 与 $2x+2y-z=2$.}
\end{opts}
\end{q}

% ---- item_id=264311e53543 seq=1.3.2.2-007 ----
\begin{q}{1.3.2.2-007}{1000第17章基础4}{ch01:q:264311e53543}
过点 $(1,0,1)$ 与点 $(0,1,1)$ 且与曲面 $z=1+x^2+y^2$ 相切的平面为 $\underline{\hspace{4em}}$。
\end{q}

% ---- item_id=c712ebc19158 seq=1.3.2.2-008 ----
\begin{q}{1.3.2.2-008}{2023数一}{ch01:q:c712ebc19158}
曲面 $z=x+2y+\ln(1+x^2+y^2)$ 在点 $(0,0,0)$ 处的切平面方程为 $\underline{\hspace{4em}}$
\end{q}

% ---- item_id=59928cafc9a2 seq=1.3.2.2-009 ----
\begin{q}{1.3.2.2-009}{1000第17章基础3}{ch01:q:59928cafc9a2}
曲面 $z=4-x^2-y^2$ 在点 $P(1,1,2)$ 处的切平面方程为 $\underline{\hspace{4em}}$.
\end{q}

% ---- item_id=44925a4f17cc seq=1.3.2.2-010 ----
\begin{q}{1.3.2.2-010}{1000第17章基础5}{ch01:q:44925a4f17cc}
已知曲面 $z=f(x,y)$ 由方程 $z-x\ln(1+z^2)+e^y=0$ 所确定，则该曲面在点 $(0,0,-1)$ 处的切平面方程为 $\underline{\hspace{4em}}$.
\end{q}

% ---- item_id=f9216f95b828 seq=1.3.2.2-011 ----
\begin{q}{1.3.2.2-011}{2024数一}{ch01:q:f9216f95b828}
已知函数 $f(x, y) = x^3 + y^3 - (x+y)^2 + 3$，设 $T$ 是曲面 $z = f(x, y)$ 在点 $(1, 1, 1)$ 处的切平面，$D$ 为 $T$ 与坐标平面所围有界区域在 $xOy$ 平面上的投影.

（\rn{I}）求 $T$ 的方程；

（\rn{II}）求 $f(x, y)$ 在 $D$ 上的最大值和最小值.
\end{q}

% ---- item_id=2df6a20c3fe0 seq=1.3.2.2-012 ----
\begin{q}{1.3.2.2-012}{1997数一}{ch01:q:2df6a20c3fe0}
（1）设直线 $l:\begin{cases}x+y+b=0,\\ x+ay-z-3=0\end{cases}$ 在平面 $\Pi$ 上，而平面 $\Pi$ 与曲面 $z=x^2+y^2$ 相切于点 $(1,-2,5)$，求 $a,b$ 的值
\end{q}

% ---- item_id=ecbbb3d2bf16 seq=1.3.2.2-013 ----
\begin{q}{1.3.2.2-013}{1989数一}{ch01:q:ecbbb3d2bf16}
(2) 已知曲面 $z=4-x^2-y^2$ 上点 $P$ 处的切平面平行于平面 $2x+2y+z-1=0$，则点 $P$ 的坐标是（ ）
\begin{opts}{0.235}
\opt{A}{$(1,-1,2)$.}\hfill\opt{B}{$(-1,1,2)$.}\hfill\opt{C}{$(1,1,2)$.}\hfill\opt{D}{$(-1,-1,2)$.}
\end{opts}
\end{q}

% ---- item_id=1b479ade0a2b seq=1.3.2.2-014 ----
\begin{q}{1.3.2.2-014}{1993数一}{ch01:q:1b479ade0a2b}
(2) 由曲线 \[\begin{cases}3x^2+2y^2=12,\\ z=0\end{cases}\] 绕 $y$ 轴旋转一周得到的旋转面在点 $(0,\sqrt{3},\sqrt{2})$ 处的指向外侧的单位法向量为 \underline{\hspace{4em}}
\end{q}

% ---- item_id=5422d937fad9 seq=1.3.2.2-015 ----
\begin{q}{1.3.2.2-015}{2000数一}{ch01:q:5422d937fad9}
曲面 $x^2+2y^2+3z^2=21$ 在点 $(1,-2,2)$ 的法线方程为 $\underline{\hspace{4em}}$
\end{q}

\par\addvspace{2.4ex}\noindent
\begin{minipage}{\linewidth}
\qhead{3}{1.3.2.3}{其他}{ch01:sec:1.3.2.3}

% ---- item_id=0ef652f7fae6 seq=1.3.2.3-001 ----
\begin{q}{1.3.2.3-001}{1000第17章基础7}{ch01:q:0ef652f7fae6}
求以 $M_0(1,1,1)$ 为顶点，以曲线 $C$（$C$ 是平面 $z=0$ 上 $y^2=x$ 被 $x=1$ 截下的有限部分）为准线的锥面方程.
\end{q}
\end{minipage}
\markboth{1.3.2.3\quad 其他}{1.3.2.3\quad 其他}


% ---- item_id=2ea8bfd6e307 seq=1.3.2.3-002 ----
\begin{q}{1.3.2.3-002}{660数一135}{ch01:q:2ea8bfd6e307}
平面 $\Pi:x-2y+2z+9=0$ 与以点 $M_0(2,0,-1)$ 为球心的球面相切，则该球面的方程是 \underline{\hspace{4em}}.
\end{q}

% ---- item_id=cd06c9598435 seq=1.3.2.3-003 ----
\begin{q}{1.3.2.3-003}{660数一344}{ch01:q:cd06c9598435}
344 曲线 $L$：\[\begin{cases}\frac{x^2}{16}+\frac{y^2}{4}-\frac{z^2}{5}=1,\\ x-2z+3=0\end{cases}\] 在 $xOy$ 面上的投影柱面方程是
\begin{opts}{0.48}
\opt{A}{$x^2+20y^2-24x-116=0.$}\hfill\opt{B}{$4y^2+4z^2-12z-7=0.$}\\[0.3ex]
\opt{C}{$\begin{cases}x^2+20y^2-24x-116=0,\\ z=0.\end{cases}$}\hfill\opt{D}{$\begin{cases}4y^2+4z^2-12z-7=0,\\ x=0.\end{cases}$}
\end{opts}
\end{q}

% ---- item_id=faded2c46c7d seq=1.3.2.3-004 ----
\begin{q}{1.3.2.3-004}{880第四章拓展}{ch01:q:faded2c46c7d}
求满足下列条件的动点的轨迹方程，并说明它们分别表示什么曲面．
（I）动点到坐标原点的距离等于它到平面 $z=4$ 的距离；
（II）动点到坐标原点的距离等于它到点 $(2,3,4)$ 的距离的一半；
（III）动点到点 $(0,0,5)$ 的距离等于它到 $x$ 轴的距离；
（IV）动点到 $x$ 轴的距离等于它到 $yOz$ 面的距离的两倍．
\end{q}

\par\addvspace{2.4ex}\noindent
\begin{minipage}{\linewidth}
\qhead{1}{1.4}{场论}{ch01:sec:1.4}

\qhead{2}{1.4.1}{方向导数}{ch01:sec:1.4.1}

\qhead{3}{1.4.1.1}{直接计算方向导数}{ch01:sec:1.4.1.1}

% ---- item_id=ddc72e8ed944 seq=1.4.1.1-001 ----
\begin{q}{1.4.1.1-001}{1991数一}{ch01:q:ddc72e8ed944}
(2) 设 $n$ 是曲面 $2x^2+3y^2+z^2=6$ 在点 $P(1,1,1)$ 处的指向外侧的法向量，求函数 $u=\frac{\sqrt{6x^2+8y^2}}{z}$ 在点 $P$ 处沿方向 $n$ 的方向导数.
\end{q}
\end{minipage}
\markboth{1.4.1.1\quad 直接计算方向导数}{1.4.1.1\quad 直接计算方向导数}


% ---- item_id=267f55726955 seq=1.4.1.1-002 ----
\begin{q}{1.4.1.1-002}{1996数一}{ch01:q:267f55726955}
函数 $u=\ln(x+\sqrt{y^2+z^2})$ 在点 $A(1,0,1)$ 处沿点 $A$ 指向点 $B(3,-2,2)$ 方向的方向导数为 $\underline{\hspace{4em}}$.
\end{q}

% ---- item_id=1f1b7855f8bc seq=1.4.1.1-003 ----
\begin{q}{1.4.1.1-003}{2005数一}{ch01:q:1f1b7855f8bc}
设函数 $u(x,y,z)=1+\frac{x^2}{6}+\frac{y^2}{12}+\frac{z^2}{18}$，单位向量 $\mathbf{n}=\frac{1}{\sqrt{3}}(1,1,1)$，则 $\left.\frac{\partial u}{\partial \mathbf{n}}\right|_{(1,2,3)} = \underline{\hspace{4em}}$
\end{q}

% ---- item_id=6d3aaf600bda seq=1.4.1.1-004 ----
\begin{q}{1.4.1.1-004}{2017数一}{ch01:q:6d3aaf600bda}
（3）函数 $f(x,y,z)=x^2y+z^2$ 在点 $(1,2,0)$ 处沿向量 $\mathbf{n}=(1,2,2)$ 的方向导数为（ ）
\begin{opts}{0.235}
\opt{A}{$12.$}\hfill\opt{B}{$6.$}\hfill\opt{C}{$4.$}\hfill\opt{D}{$2.$}
\end{opts}
\end{q}

% ---- item_id=6858e38ba5bc seq=1.4.1.1-005 ----
\begin{q}{1.4.1.1-005}{2025数一}{ch01:q:6858e38ba5bc}
已知函数 $v(x,y,z)=xy^2z^3$，向量 $\vec{n}=(2,2,-1)$，则 $\left.\frac{\partial v}{\partial \vec{n}}\right|_{(1,1,1)}=\underline{\hspace{4em}}$。
\end{q}

% ---- item_id=803519ac016d seq=1.4.1.1-006 ----
\begin{q}{1.4.1.1-006}{1000第17章基础6}{ch01:q:803519ac016d}
设 $z=\ln\left(\mathrm{e}^{-y}+\dfrac{y^2}{x}\right)$ 在点 $(1,1)$ 处沿 $l=(1,0)$ 的方向导数为 \underline{\hspace{4em}}．
\end{q}

% ---- item_id=1d74aa249fa2 seq=1.4.1.1-007 ----
\begin{q}{1.4.1.1-007}{1000第17章基础8}{ch01:q:1d74aa249fa2}
设函数 $f(x,y)=\begin{cases}\dfrac{x^2|y|}{x^2+y^2}, & (x,y)\neq(0,0),\\ 0, & (x,y)=(0,0),\end{cases}$，则 $f(x,y)$ 在点 $(0,0)$ 处沿 $l=(1,1)$ 的方向导数是 \underline{\hspace{4em}}．
\end{q}

% ---- item_id=e1271ff139d7 seq=1.4.1.1-008 ----
\begin{q}{1.4.1.1-008}{660数一360}{ch01:q:e1271ff139d7}
函数 $f(x,y,z)=x^2+y^2+z^2$ 在点 $(1,-1,1)$ 处沿曲线 $x=t,y=-t^2,z=t^3$ 在该点指向 $x$ 轴负向一侧的切线方向的方向导数等于
\begin{opts}{0.235}
\opt{A}{$-12$.}\hfill\opt{B}{$12$.}\hfill\opt{C}{$-\frac{12}{\sqrt{14}}$.}\hfill\opt{D}{$\frac{12}{\sqrt{14}}$.}
\end{opts}
\end{q}

% ---- item_id=78d68f397aa4 seq=1.4.1.1-009 ----
\begin{q}{1.4.1.1-009}{1000第17章基础13}{ch01:q:78d68f397aa4}
设可微函数 $z=z(x,y)$ 在平面上任一点 $(x,y)$ 处沿 $x$ 轴正向 $i$ 与 $y$ 轴正向 $j$ 的方向导数分别为 $[e^{-x}-f(x)]y$ 与 $f(x)$，其中 $f(x)$ 的一阶导数连续，且 $f(0)=1$.
(1) 求 $z(x,y)$ 的表达式；
(2) 判断 $z(x,y)$ 是否有极值，若有，求之，若无，说明理由.
\end{q}

\par\addvspace{2.4ex}\noindent
\begin{minipage}{\linewidth}
\qhead{3}{1.4.1.2}{最大/最小方向导数}{ch01:sec:1.4.1.2}

% ---- item_id=30fa9cacfaba seq=1.4.1.2-001 ----
\begin{q}{1.4.1.2-001}{2002数一}{ch01:q:30fa9cacfaba}
设有一小山，取它的底面所在的平面为 $xOy$ 坐标面，其底部所占的区域为
\[
D=\{(x,y)\mid x^2+y^2-xy\leqslant 75\},
\]
小山高度函数为
\[
h(x,y)=75-x^2-y^2+xy.
\]
（I）设 $M(x_0,y_0)$ 为区域 $D$ 上一点，问 $h(x,y)$ 在该点沿平面上什么方向的方向导数最大？若记此方向导数的最大值为 $g(x_0,y_0)$，试写出 $g(x_0,y_0)$ 的表达式；
（II）现欲利用此小山开展攀岩活动，为此需要在山脚寻找一上山坡度最大的点作为攀登的起点。也就是说，要在 $D$ 的边界线 $x^2+y^2-xy=75$ 上找出使（I）中的 $g(x,y)$ 达到最大值的点。试确定攀登起点的位置。
\end{q}
\end{minipage}
\markboth{1.4.1.2\quad 最大/最小方向导数}{1.4.1.2\quad 最大/最小方向导数}


% ---- item_id=26f9c73f068d seq=1.4.1.2-002 ----
\begin{q}{1.4.1.2-002}{2015数一}{ch01:q:26f9c73f068d}
已知函数 $f(x,y)=x+y+xy$，曲线 $C:x^{2}+y^{2}+xy=3$，求 $f(x,y)$ 在曲线 $C$ 上的最大方向导数.
\end{q}

% ---- item_id=5cab364221e9 seq=1.4.1.2-003 ----
\begin{q}{1.4.1.2-003}{2022数一}{ch01:q:5cab364221e9}
函数 $f(x,y)=x^2+2y^2$ 在点 $(0,1)$ 处的最大方向导数为 $\underline{\hspace{4em}}$
\end{q}

% ---- item_id=f78c16b34703 seq=1.4.1.2-004 ----
\begin{q}{1.4.1.2-004}{1000第17章强化8}{ch01:q:f78c16b34703}
设可微函数 $f(u,v)$ 满足 $f(x-y,x+e^y)=x^2-y^2$，则 $f(u,v)$ 在点 $(1,2)$ 处的方向导数的最大值等于（  ）．
\begin{opts}{0.235}
\opt{A}{$1$}\hfill\opt{B}{$\sqrt{2}$}\hfill\opt{C}{$\sqrt{3}$}\hfill\opt{D}{$2$}
\end{opts}
\end{q}

% ---- item_id=ff2602660288 seq=1.4.1.2-005 ----
\begin{q}{1.4.1.2-005}{1000第17章强化10}{ch01:q:ff2602660288}
已知函数 $z=f(x,y)$ 可微，其在点 $P_0(1,2)$ 处沿从 $P_0$ 到 $P_1(2,3)$ 的方向的方向导数为 $2\sqrt{2}$，沿从 $P_0$ 到 $P_2(1,0)$ 的方向的方向导数为 $-3$，则 $z$ 在点 $P_0$ 处的最大方向导数为 $\underline{\hspace{4em}}$。
\end{q}

% ---- item_id=25ace12b250c seq=1.4.1.2-006 ----
\begin{q}{1.4.1.2-006}{1000第17章强化11}{ch01:q:25ace12b250c}
函数 $u(x,y,z)=xy-2z^2$ 在点 $(1,1,-2)$ 处的最大方向导数为 \underline{\hspace{4em}}.
\end{q}

% ---- item_id=6eded9bf0338 seq=1.4.1.2-007 ----
\begin{q}{1.4.1.2-007}{660数一350}{ch01:q:6eded9bf0338}
函数 $f(x,y,z)=x^2y^3+3y^2z^3$ 在点 $(0,1,1)$ 处方向导数的最大值为
\begin{opts}{0.235}
\opt{A}{$\sqrt{107}.$}\hfill\opt{B}{$\sqrt{117}.$}\hfill\opt{C}{$117.$}\hfill\opt{D}{$107.$}
\end{opts}
\end{q}

% ---- item_id=c9d8c08835c4 seq=1.4.1.2-008 ----
\begin{q}{1.4.1.2-008}{1000第17章基础9}{ch01:q:c9d8c08835c4}
函数 $f(x,y)=2x^2+y^2$ 在点 $(0,1)$ 处的最大方向导数为 $\underline{\hspace{4em}}$.
\end{q}

% ---- item_id=3416e2eb7720 seq=1.4.1.2-009 ----
\begin{q}{1.4.1.2-009}{1000第17章基础12}{ch01:q:3416e2eb7720}
设 $g(x,y)$ 是函数 $f(x,y)=x+2y+xy$ 在点 $(x,y)$ 处的最大方向导数.
(1) 求 $g(x,y)$ 的表达式；
(2) 求 $g(x,y)$ 在曲线 $C:x^2+y^2=5$ 上的最大值.
\end{q}

% ---- item_id=b14e43c35093 seq=1.4.1.2-010 ----
\begin{q}{1.4.1.2-010}{1000第17章强化14}{ch01:q:b14e43c35093}
在曲面 $\Sigma:x^2+2y^2+z^2=1$ 上求一点，使函数 $f(x,y,z)=2x^2+y^2+z$ 在该点沿方向 $n$ 的方向导数最大，其中 $n$ 是曲面 $\Sigma$ 在点 $P\left(\frac{1}{2},-\frac{1}{2},\frac{1}{2}\right)$ 的外侧法向量.
\end{q}

% ---- item_id=dc97102e918e seq=1.4.1.2-011 ----
\begin{q}{1.4.1.2-011}{1000第17章基础10}{ch01:q:dc97102e918e}
设 $f(x,y,z)=x^2 e^{y z^2}$，则 $f(x,y,z)$ 在点 $(-1,0,1)$ 处的方向导数的最小值为 \underline{\hspace{4em}}。
\end{q}

\par\addvspace{2.4ex}\noindent
\begin{minipage}{\linewidth}
\qhead{3}{1.4.1.3}{与可微综合}{ch01:sec:1.4.1.3}

% ---- item_id=1e6a3b97a4ce seq=1.4.1.3-001 ----
\begin{q}{1.4.1.3-001}{2020数一}{ch01:q:1e6a3b97a4ce}
(3) 设函数 $f(x,y)$ 在点 $(0,0)$ 处可微，$f(0,0)=0$，$\mathbf{n}=\left.\left(\frac{\partial f}{\partial x},\frac{\partial f}{\partial y},-1\right)\right|_{(0,0)}$，非零向量 $\boldsymbol{\alpha}$ 与 $\mathbf{n}$ 垂直，则（　　）.
\begin{opts}{0.48}
\opt{A}{$\displaystyle \lim\limits_{(x,y)\to(0,0)}\dfrac{|\mathbf{n}\cdot(x,y,f(x,y))|}{\sqrt{x^{2}+y^{2}}}$ 存在}\hfill\opt{B}{$\displaystyle \lim\limits_{(x,y)\to(0,0)}\dfrac{|\mathbf{n}\times(x,y,f(x,y))|}{\sqrt{x^{2}+y^{2}}}$ 存在}\\[0.3ex]
\opt{C}{$\displaystyle \lim\limits_{(x,y)\to(0,0)}\dfrac{|\boldsymbol{\alpha}\cdot(x,y,f(x,y))|}{\sqrt{x^{2}+y^{2}}}$ 存在}\hfill\opt{D}{$\displaystyle \lim\limits_{(x,y)\to(0,0)}\dfrac{|\boldsymbol{\alpha}\times(x,y,f(x,y))|}{\sqrt{x^{2}+y^{2}}}$ 存在}
\end{opts}
\end{q}
\end{minipage}
\markboth{1.4.1.3\quad 与可微综合}{1.4.1.3\quad 与可微综合}


% ---- item_id=2fa151ab7257 seq=1.4.1.3-002 ----
\begin{q}{1.4.1.3-002}{1000第17章基础1}{ch01:q:2fa151ab7257}
已知函数 $f(x,y)$ 在点 $(0,0)$ 处可微，$f(0,0)=0$，$f'_x(0,0)=1$，$f'_y(0,0)=-1$，且 $\boldsymbol{n}=(-1,1,1)$，则
\[\lim_{\substack{x\to 0\\ y\to 0}} \frac{(x,y,f(x,y))\cdot \boldsymbol{n}}{e^{\sqrt{x^2+y^2}}-1} = \underline{\hspace{4em}} .\]
\end{q}

\par\addvspace{2.4ex}\noindent
\begin{minipage}{\linewidth}
\qhead{2}{1.4.2}{梯度}{ch01:sec:1.4.2}

% ---- item_id=5f5fb5b2ed14 seq=1.4.2-001 ----
\begin{q}{1.4.2-001}{1992数一}{ch01:q:5f5fb5b2ed14}
函数 $u=\ln(x^2+y^2+z^2)$ 在点 $M(1,2,-2)$ 处的梯度 $\left.\operatorname{grad}u\right|_M=\underline{\hspace{4em}}$
\end{q}
\end{minipage}
\markboth{1.4.2\quad 梯度}{1.4.2\quad 梯度}


% ---- item_id=243891b20bb1 seq=1.4.2-002 ----
\begin{q}{1.4.2-002}{2008数一}{ch01:q:243891b20bb1}
(2) 函数 $f(x,y)=\arctan \frac{x}{y}$ 在点 $(0,1)$ 处的梯度等于（ ）
\begin{opts}{0.235}
\opt{A}{$\boldsymbol{i}$．}\hfill\opt{B}{$-\boldsymbol{i}$．}\hfill\opt{C}{$\boldsymbol{j}$．}\hfill\opt{D}{$-\boldsymbol{j}$．}
\end{opts}
\end{q}

% ---- item_id=8183e6f34cc7 seq=1.4.2-003 ----
\begin{q}{1.4.2-003}{2012数一}{ch01:q:8183e6f34cc7}
$\displaystyle \operatorname{grad}\left(xy+\frac{z}{y}\right)\bigg|_{(2,1,1)} = \underline{\hspace{4em}}$
\end{q}

% ---- item_id=f1397f87d548 seq=1.4.2-004 ----
\begin{q}{1.4.2-004}{1000第17章强化9}{ch01:q:f1397f87d548}
$\displaystyle f(x,y,z)=x^2\int_1^y e^t\,dt+2z$ 在点 $(1,1,2)$ 处的梯度为 $\underline{\hspace{4em}}$.
\end{q}

% ---- item_id=2bc2a1a11e08 seq=1.4.2-005 ----
\begin{q}{1.4.2-005}{1000第17章强化12}{ch01:q:2bc2a1a11e08}
设可微函数 $f(x,y)$ 在点 $(0,0)$ 处的最小方向导数为 $a, a\neq 0$ 且 $a\neq -1, b,c$ 是满足 $b^2+c^2$ 为正常数的任意实数, 则 $\operatorname{grad} f(0,0)$ 与 $(b,c)$ 内积的最大值为 ( ).
\begin{opts}{0.235}
\opt{A}{$a\sqrt{b^2+c^2}$}\hfill\opt{B}{$-a\sqrt{b^2+c^2}$}\hfill\opt{C}{$\sqrt{|a|}(b^2+c^2)$}\hfill\opt{D}{$-\sqrt{|a|}(b^2+c^2)$}
\end{opts}
\end{q}

% ---- item_id=ae39c9075451 seq=1.4.2-006 ----
\begin{q}{1.4.2-006}{1000第17章强化13}{ch01:q:ae39c9075451}
设 $f(x,y)$ 可微, $p(x,y), p_0(x_0,y_0)$ 为曲面 $z=f(x,y)$ 上的点, 则 ( ).
\begin{opts}{0.97}
\opt{A}{$\displaystyle \lim_{p\to p_0} \frac{f(p)+\left[f(p_0)-\operatorname{grad}f\big|_{p_0}\circ \overrightarrow{p_0 p}\right]}{\left|\overrightarrow{p_0 p}\right|}=0$}\\[0.3ex]
\opt{B}{$\displaystyle \lim_{p\to p_0} \frac{f(p)-\left[f(p_0)-\operatorname{grad}f\big|_{p_0}\cdot \overrightarrow{p_0 p}\right]}{\left|\overrightarrow{p_0 p}\right|}=0$}\\[0.3ex]
\opt{C}{$\displaystyle \lim_{p\to p_0} \frac{f(p)+\left[f(p_0)+\operatorname{grad}f\big|_{p_0}\cdot \overrightarrow{p_0 p}\right]}{\left|\overrightarrow{p_0 p}\right|}=0$}\\[0.3ex]
\opt{D}{$\displaystyle \lim_{p\to p_0} \frac{f(p)-\left[f(p_0)+\operatorname{grad}f\big|_{p_0}\cdot \overrightarrow{p_0 p}^2\right]}{\left|\overrightarrow{p_0 p}\right|}=0$}
\end{opts}
\end{q}

% ---- item_id=1d505a0fad35 seq=1.4.2-007 ----
\begin{q}{1.4.2-007}{1000第17章强化1}{ch01:q:1d505a0fad35}
设 $f(x,y)=\mathrm{e}^{-(x^2+2y^2)}$，曲线 $y=y(x)$ 上任一点 $P(x,y)$ 的切线方向始终指向 $f(x,y)$ 变化率最大的方向，且 $y(1)=2$，则 $y(x)=(\quad)$。
\begin{opts}{0.235}
\opt{A}{$2x^2$}\hfill\opt{B}{$x^2+x$}\hfill\opt{C}{$e^{-x^2+1}+x$}\hfill\opt{D}{$2e^{-x^2+1}$}
\end{opts}
\end{q}

\par\addvspace{2.4ex}\noindent
\begin{minipage}{\linewidth}
\qhead{2}{1.4.3}{散度}{ch01:sec:1.4.3}

% ---- item_id=6dfb16461bbe seq=1.4.3-001 ----
\begin{q}{1.4.3-001}{1989数一}{ch01:q:6dfb16461bbe}
向量场 $\mathbf{u}(x,y,z)=xy^{2}\mathbf{i}+ye^{z}\mathbf{j}+x\ln(1+z^{2})\mathbf{k}$ 在点 $P(1,1,0)$ 处的散度 $\operatorname{div}\mathbf{u}=\underline{\hspace{4em}}$
\end{q}
\end{minipage}
\markboth{1.4.3\quad 散度}{1.4.3\quad 散度}


% ---- item_id=221dfbdd54eb seq=1.4.3-002 ----
\begin{q}{1.4.3-002}{2026数一}{ch01:q:221dfbdd54eb}
设向量 $\boldsymbol{v}_1=(0,x,z)$，$\boldsymbol{v}_2=(y,0,1)$，令 $\boldsymbol{F}(x,y,z)=\boldsymbol{v}_1\times \boldsymbol{v}_2$，则 $\operatorname{div}\boldsymbol{F}=\underline{\hspace{4em}}$。
\end{q}

% ---- item_id=a76185246213 seq=1.4.3-003 ----
\begin{q}{1.4.3-003}{900第九章A类16}{ch01:q:a76185246213}
设向量场 $\boldsymbol{F}(x,y,z)=\frac{x}{(x^2+2y^2+3z^2)^2}\boldsymbol{i}+\frac{y}{(x^2+2y^2+3z^2)^2}\boldsymbol{j}+\frac{z}{(x^2+2y^2+3z^2)^2}\boldsymbol{k}$，则 $|\operatorname{div}\boldsymbol{F}|$ 在单位球面 $x^2+y^2+z^2=1$ 上的最大值为 $\underline{\hspace{4em}}$。
\end{q}

\par\addvspace{2.4ex}\noindent
\begin{minipage}{\linewidth}
\qhead{2}{1.4.4}{旋度}{ch01:sec:1.4.4}

% ---- item_id=2e4ba9cd540c seq=1.4.4-001 ----
\begin{q}{1.4.4-001}{2016数一; 660数一151}{ch01:q:2e4ba9cd540c}
向量场 $\mathbf{A}(x,y,z)=(x+y+z)\mathbf{i}+xy\mathbf{j}+z\mathbf{k}$ 的旋度 $\operatorname{rot}\mathbf{A}=\underline{\hspace{4em}}$
\end{q}
\end{minipage}
\markboth{1.4.4\quad 旋度}{1.4.4\quad 旋度}


% ---- item_id=1a6cc688a5d8 seq=1.4.4-002 ----
\begin{q}{1.4.4-002}{2018数一}{ch01:q:1a6cc688a5d8}
设 $\mathbf{F}(x,y,z)=xy\mathbf{i}-yz\mathbf{j}+zx\mathbf{k}$，则 $\operatorname{rot}\mathbf{F}(1,1,0)=\underline{\hspace{4em}}$
\end{q}

% ---- item_id=4b955c4f1a2e seq=1.4.4-003 ----
\begin{q}{1.4.4-003}{1000第17章基础14}{ch01:q:4b955c4f1a2e}
设 $F(x,y,z)=xyi-y\cos zj+z\sin xk$，则 $\operatorname{rot}F\big|_{(1,1,0)}=\underline{\hspace{4em}}$.
\end{q}

% ---- item_id=ca7c14657360 seq=1.4.4-004 ----
\begin{q}{1.4.4-004}{1000第17章强化15}{ch01:q:ca7c14657360}
设 $ \boldsymbol{F}(x,y,z)=|x|\cos y\,\boldsymbol{i}-y\sin z\,\boldsymbol{j}+z\boldsymbol{k} $，则 $ \operatorname{rot}\boldsymbol{F}(1,1,0)=\underline{\hspace{4em}} $。
\end{q}

% ---- item_id=f8d2b600566e seq=1.4.4-005 ----
\begin{q}{1.4.4-005}{900第九章A类17}{ch01:q:f8d2b600566e}
向量场 $\boldsymbol{u}(x,y,z)=xy^2\boldsymbol{i}+yz^2\boldsymbol{j}+zx^2\boldsymbol{k}$ 在点 $M(1,1,1)$ 处的旋度 $\operatorname{rot}\boldsymbol{u}=\underline{\hspace{4em}}$。
\end{q}

% ---- item_id=475b0d1a029f seq=1.4.4-006 ----
\begin{q}{1.4.4-006}{900第九章A类18}{ch01:q:475b0d1a029f}
设向量场 $\boldsymbol{A}(x,y,z)=e^{xy}\boldsymbol{i}+y\ln z\,\boldsymbol{j}+\sin xyz\,\boldsymbol{k}$，则其在点 $(0,1,1)$ 处的旋度 $\operatorname{rot}\boldsymbol{A}\big|_{(0,1,1)}$ 的长度为 $\underline{\hspace{4em}}$．
\end{q}

\par\addvspace{2.4ex}\noindent
\begin{minipage}{\linewidth}
\qhead{2}{1.4.5}{混合}{ch01:sec:1.4.5}

% ---- item_id=6fa9ab2d7642 seq=1.4.5-001 ----
\begin{q}{1.4.5-001}{1993数一}{ch01:q:6fa9ab2d7642}
设数量场 $u=\ln\sqrt{x^2+y^2+z^2}$，则 $\operatorname{div}(\operatorname{grad}u)=\underline{\hspace{4em}}$
\end{q}
\end{minipage}
\markboth{1.4.5\quad 混合}{1.4.5\quad 混合}


% ---- item_id=ede3db9c06d7 seq=1.4.5-002 ----
\begin{q}{1.4.5-002}{2001数一; 660数一160}{ch01:q:ede3db9c06d7}
设 $r=\sqrt{x^2+y^2+z^2}$，则 $\mathrm{div}(\mathbf{grad}\,r)\big|_{(1,-2,2)}=\underline{\hspace{4em}}$
\end{q}

% ---- item_id=bb2c5880a33d seq=1.4.5-003 ----
\begin{q}{1.4.5-003}{880第九章综合填空6}{ch01:q:bb2c5880a33d}
(6) 向量场 $A(x,y,z)=(x+y+z)\boldsymbol{i}+xy\boldsymbol{j}+z\boldsymbol{k}$ 在点 $P(1,1,1)$ 处的旋度 $\mathrm{rot}A=\underline{\hspace{4em}}$，$\mathrm{div}A=\underline{\hspace{4em}}$。
\end{q}
